% ************************** Thesis Abstract *****************************
% Use `abstract' as an option in the document class to print only the titlepage and the abstract.
\begin{abstract}
Understanding visual scenes at a holistic and structured level has long been a central goal in computer vision. Among various formulations, Panoptic Scene Graph Generation (PSG) stands out as a comprehensive task that simultaneously segments all entities in an image and infers their pairwise relations, thereby constructing a semantically rich, pixel-level scene graph. Such representations are crucial for downstream applications including visual question answering, image captioning, embodied navigation, and robotic reasoning. Despite recent progress, existing PSG approaches suffer from three key limitations: (1) the long-tail imbalance in relation categories, which biases models toward common relations and suppresses rare but semantically valuable ones; (2) the lack of multimodal reasoning, as current frameworks rely solely on visual inputs and fail to exploit linguistic or commonsense knowledge; and (3) the closed-set assumption that constrains recognition to predefined object and relation categories, hindering scalability in open-world scenarios.
This thesis systematically addresses these challenges through three successive frameworks: HiLo, VLPrompt, and OpenPSG, which together trace an evolution from bias mitigation, to vision-language integration, and finally to open-set generalization.

The first part of this thesis introduces HiLo, the first explicitly unbiased framework for Panoptic Scene Graph Generation. Unlike prior re-sampling or class-balanced loss methods that treat the long-tail issue as a static data imbalance problem, HiLo recognizes the phenomenon of relational semantic overlap: subject–object pairs often share multiple overlapping relations spanning both high- and low-frequency categories (e.g., ``beside'', ``playing'', ``chasing''). To address this, HiLo employs two parallel network branches specialized for high- and low-frequency relations, respectively. These branches are trained jointly with novel subject–object consistency and relation consistency losses, which ensure cross-branch alignment while allowing specialization. A transformer-based baseline derived from Mask2Former serves as the architectural foundation, enabling dense, panoptic-level reasoning. The resulting model significantly improves both global and mean recall on benchmark datasets, setting a new standard for unbiased PSG performance and offering a principled treatment of relational imbalance at the feature-learning level.

Building upon this foundation, the second part of the thesis extends the scope of PSG beyond vision-only reasoning by introducing VLPrompt, a Vision-Language Prompting framework that integrates the powerful commonsense reasoning capabilities of Large Language Models (LLMs) into relation prediction. While prior work in scene graph generation occasionally exploited static language priors such as relation names or knowledge graphs, VLPrompt harnesses dynamically generated, context-rich linguistic information through prompt engineering and chain-of-thought reasoning. Specifically, two forms of prompts—relation proposer and relation judger—elicit LLM-generated textual descriptions about the plausibility and semantics of relations between subject–object pairs. These descriptions are encoded into language features and fused with visual embeddings via a dual-attention vision-language prompter network. This design enables complementary learning: the proposer captures frequent, generic relations, while the judger refines rare or context-dependent ones. Extensive experiments on the PSG dataset demonstrate that VLPrompt dramatically boosts both recall and mean recall (e.g., R@100 +9.4, mR@100 +20.6 over prior state-of-the-art), validating the effectiveness of language-augmented relational reasoning and establishing the first LLM-driven model for the PSG task.

The third contribution, OpenPSG, moves beyond the closed-world assumption by enabling open-set panoptic scene graph generation. Leveraging Large Multimodal Models (LMMs) such as BLIP-2, OpenPSG performs open-set relation prediction in an autoregressive manner. It introduces a Relation Query Transformer that efficiently extracts subject–object pair features and predicts the likelihood of relational existence, filtering irrelevant pairs to improve inference efficiency. A Multimodal Relation Decoder, guided by generation and judgment instructions, then utilizes the LMM’s generative capacity to produce or verify relation descriptions for both known and novel predicates. Combined with an open-set panoptic segmenter (e.g., OpenSeeD), this design enables the joint recognition of unseen object categories and novel relations, achieving structured understanding beyond predefined label vocabularies. OpenPSG thus represents the first comprehensive framework to unify open-set segmentation and open-set relational reasoning, bridging the gap between structured scene understanding and large-scale multimodal intelligence.

Together, HiLo, VLPrompt, and OpenPSG delineate a coherent research trajectory: from explicitly unbiased relation learning, to multimodal reasoning through language prompting, to open-world generalization powered by large multimodal models. Collectively, these works redefine the landscape of Panoptic Scene Graph Generation by progressively expanding its representational, cognitive, and generalization capacities. The thesis not only advances PSG as a research field but also provides a conceptual and methodological foundation for future developments in multimodal, interpretable, and scalable visual reasoning systems capable of understanding complex scenes in the open world.
\end{abstract}
